\documentclass[12pt]{article}
\usepackage[utf8]{inputenc}
\usepackage[T1]{fontenc}
\usepackage{hyperref}
\title{Tapachultec Language Resources for K-12 Curriculum Development}
\author{Compiled by Claude Sonnet 5.5 (claude-sonnet-5-5) for Andr\'as Kornai}
\date{}
\begin{document}
\maketitle

\section{Introduction}
Tapachultec (ISO 639-3 qcs; Glottocode tapa1260) is an extinct Mixean language of the Mixe--Zoque family, once spoken around Tapachula in Chiapas, Mexico. Sources date its extinction to about the 1930s. It was documented only fragmentarily, so no community, school programme or native-language teaching material was found in the searches. A K-12 curriculum in the language itself is not possible with current documentation.

\section{Primary Educational Resources}
No primers, textbooks or school materials were found.

\section{Linguistic and Reference Materials}
\begin{itemize}
\item Wikipedia, Tapachultec language: \url{https://en.wikipedia.org/wiki/Tapachultec_language}
\item Wikipedia, Mixe--Zoque languages: \url{https://en.wikipedia.org/wiki/Mixe%E2%80%93Zoque_languages}
\item Wikipedia, Mixean languages: \url{https://en.wikipedia.org/wiki/Mixean_languages}
\item Britannica entry on Tapachultec: \url{https://www.britannica.com/topic/Tapachultec}
\item Glottolog reference, Terrence S. Kaufman 1964 (on the position of Tapachulteco within Mixe--Zoque): \url{https://glottolog.org/resource/reference/id/128593}
\item Glottolog reference, Pablo Gonz\'alez Casanova 1927 (scanty lexical materials): \url{https://glottolog.org/resource/reference/id/320800}
\item ASJP wordlist page: \url{https://asjp.clld.org/languages/TAPACHULTEC}
\end{itemize}

\section{Religious and Community Resources}
None found. The language has no known speakers or community organisation.

\section{Digital and Audio Resources}
No audio or software was found. A map entry exists at Native Land: \url{https://native-land.ca/maps/languages/tapachulteco-mixe-zoque-2}

\section{Curriculum Development Context}
Because the language is extinct and the record is thin (mostly a 16th-century description and early 20th-century word lists), any curriculum could only be a heritage-awareness unit about the language and its place in the Mixe--Zoque family, taught in Spanish or English. Closely related living languages would be better bases for immersion work.
\end{document}